% Part: proof-theory
% Chapter: sequent-calculus
% Section: quantifiers
%READERNOTE{SOL6-A313: क्रमवर्ती संख्यापक-नियमांच्या बंद पदे वापरण्याच्या संकेताशी आदेशनाची अट जुळवली आहे; आधीची नियमितीकरण आणि विगमनातील घट जपली आहे.}

\documentclass[../../../include/open-logic-section]{subfiles}

\begin{document}

\olfileid{pt}{seq}{qua}

\olsection{नियमित सिद्धता आणि प्रतिस्थापन}

संख्यकांचे नियम काही गुंतागुंत निर्माण करतात;
कारण~\LeftR{\lexists} आणि~\RightR{\lforall} या
नियमांना ``आयगेन चराच्या अटी'' लागू होतात. या
बहुतेक अडचणी त्या निगमनांतील !!{constant}~$a$
पुन्हा कधीही वापरला जाणार नाही याची खात्री करून
सोडवता येतात. म्हणून पुढील व्याख्या देऊ.

\begin{defn}
\Log{G1c} किंवा \Log{G3c} मधील !!^a{proof}~$\pi$
\emph{नियमित} असते, जर
\begin{enumerate}
  \item कोणताही आयगेन चर~$a$ एकाहून अधिक
  \LeftR{\lexists} किंवा~\RightR{\lforall} निगमनांचा
  आयगेन चर म्हणून वापरला जात नसेल; आणि
  \item $a$ हा \LeftR{\lexists} किंवा~\RightR{\lforall}
  निगमनाचा आयगेन चर असेल, तर तो~$\pi$ मध्ये
  फक्त त्या निगमनाच्या वरच येत असेल.
\end{enumerate}
\end{defn}

\begin{lem}\ollabel{lem:subst-regular} $\pi$ ही
नियमित असून तिची अंतिम क्रमवर्ती $\Gamma \Sequent
\Delta$ आहे असे माना. तसेच~$c$ हा~$\pi$ मध्ये
आयगेन चर म्हणून न वापरलेला !!a{constant} असो,
आणि~$t$ हे~$\pi$ चा कोणताही आयगेन चर नसलेले बंद पद असो.
येथील संख्यापक-नियमांत पद बंद असण्याची अट आहे;
त्या संकेतानुसार ही आदेशनाची मर्यादा आवश्यक आहे.
मग~$\pi$ मधील~$c$ ची प्रत्येक उपस्थिती~$t$ ने
बदलून मिळणारी~$\Subst{\pi}{t}{c}$ ही नियमित
!!{proof} असते. $\Subst{\pi}{t}{c}$ ची अंतिम क्रमवर्ती
$\Subst{\Gamma}{t}{c}
\Sequent \Subst{\Delta}{t}{c}$ असते; येथे
$\Subst{\Gamma}{t}{c} =
\Setabs{!A(t)}{!A(c) \in \Gamma}$.
\end{lem}

\begin{proof}
  $\pheight{\pi}$ वर विगमन करू. $\pheight{\pi} = 0$
  असेल तर ती फक्त एक स्वयंसिद्धक आहे. मग
  $\Subst{\pi}{t}{c}$ हेही स्वयंसिद्धक असते, कारण
  त्यातील कोणतेही !!{formula}~$!A(c)$ हे~$!A(t)$
  होते.

$\pheight{\pi} > 0$ असेल तर~$\pi$
  मधील शेवटच्या निगमनानुसार प्रसंग वेगळे करू.
  संरचनात्मक नियमांचे (\Log{G1c} मध्ये) आणि
  विधानात्मक संयोजकांच्या तार्किक नियमांचे प्रसंग
  सोपे आहेत; ते प्रश्न म्हणून सोडतो. $\pi$ चा शेवट
  संख्यकाच्या निगमनाने होणारे प्रसंग पाहू.
  \begin{enumerate}
    \item शेवटचे निगमन $\RightR{\lforall}$ असेल, तर
    $\pi$ या रूपाची असते:
\[
    \AxiomC{}
    \RightLabel{$\pi'$}
    \Deduce$\Gamma(c) \fCenter \Delta'(c), !A(a,c)$
    \RightLabel{\RightR{\lforall}}
    \UnaryInf$\Gamma(c) \fCenter \Delta'(c), \lforall[x][!A(x,c)]$
    \DisplayProof
    \]
    विगमनात्मक गृहीतकानुसार $\Subst{\pi'}{t}{c}$
    ही $\Gamma(t) \fCenter \Delta'(t), !A(a,t)$ ची
    नियमित !!{proof} आहे. $a$ हा~$\pi$ चा आयगेन
    चर असल्यामुळे~$a$ हे~$t$ मध्ये येत नाही.
    म्हणून $\Subst{\pi}{t}{c}$ मधील शेवटचे निगमन,
\[
    \AxiomC{}
    \RightLabel{$\Subst{\pi'}{t}{c}$}
    \Deduce$\Gamma(t) \fCenter \Delta'(t), !A(a,t)$
    \RightLabel{\RightR{\lforall}}
    \UnaryInf$\Gamma(t) \fCenter \Delta'(t), \lforall[x][!A(x,t)]$
    \DisplayProof
    \]
    योग्य आहे: त्याच्या निष्कर्षात~$a$ येत नाही.
    \item शेवटचे निगमन \Log{G3c} मधील
    $\LeftR{\lforall}$ असेल, तर $\pi$ चे रूप असे:
\[
    \AxiomC{}
    \RightLabel{$\pi'$}
    \Deduce$!A(s(c),c), \lforall[x][!A(x,c)], \Gamma(c) \fCenter \Delta'(c)$
    \RightLabel{\LeftR{\lforall}}
    \UnaryInf$\lforall[x][!A(x,c)], \Gamma(c) \fCenter \Delta'(c)$
    \DisplayProof
    \]
    विगमनात्मक गृहीतकानुसार $\Subst{\pi'}{t}{c}$
    ही $!A(s(t),t), \lforall[x][!A(x,t)], \Gamma(t) \fCenter
    \Delta'(t)$ ची नियमित !!{proof} आहे.
    $\Subst{\pi}{t}{c}$ मधील शेवटचे निगमन,
\[
    \AxiomC{}
    \RightLabel{$\Subst{\pi'}{t}{c}$}
    \Deduce$!A(s(t),t), \lforall[x][!A(x,t)], \Gamma(t) \fCenter \Delta'(t)$
    \RightLabel{\LeftR{\lforall}}
    \UnaryInf$\lforall[x][!A(x,t)], \Gamma(t) \fCenter \Delta'(t)$
    \DisplayProof
    \]
    योग्य आहे. \Log{G1c} मधील \LeftR{\lforall}
    साठीचा प्रसंगही असाच, पण काहीसा सोपा आहे.
    \item शेवटचे निगमन~\RightR{\lexists} किंवा
    \LeftR{\lexists} असेल, तर ते प्रश्न म्हणून सोडले आहे.
  \end{enumerate}
  प्रत्येक प्रसंगात आपण नियमित !!{proof} च्या
  शेवटी एक क्रमवर्ती जोडली (दोन पूर्वपक्ष असलेल्या
  नियमात दोन नियमित !!{proof}s वापरल्या). नव्या
  अंतिम क्रमवर्तीमध्ये~$\pi$ च्या अंतिम
  क्रमवर्तीपेक्षा फक्त~$c$ च्या जागी पद~$t$ आलेले
  असते. तसेच $\pi$ आणि $\Subst{\pi}{t}{c}$
  यांचे आयगेन चर तेच राहतात; कारण~$t$ मध्ये
  $\pi$ चा कोणताही आयगेन चर नाही, नव्या अंतिम
  क्रमवर्तीमध्येही $\Subst{\pi}{t}{c}$ चा आयगेन चर
  येत नाही. त्यामुळे~$\Subst{\pi}{t}{c}$ नियमित असते.
\end{proof}

\begin{prop}
$\pi$ ही \Log{G1c} किंवा \Log{G3c} मधील
!!a{proof} असेल, तर तिच्याशी तीच रचना
(विशेषतः तीच !!{height}) असलेली नियमित
!!{proof}~$\pi'$ मिळते, जी फक्त आयगेन चर
बदलल्यामुळे~$\pi$ पेक्षा वेगळी असते.
\end{prop}

\begin{proof}
फक्त या सिद्धतेपुरते,~$\pi$ मधील आयगेन चराचे
एखादे निगमन \emph{स्वच्छ} म्हणू, जर त्याचा आयगेन
चर इतर कोणत्याही निगमनाचा आयगेन चर नसेल आणि
त्या निगमनाच्या तत्काळ उप-!!{proof} व्यतिरिक्त
~$\pi$ मध्ये कुठेही येत नसेल; अन्यथा ते
\emph{अस्वच्छ} म्हणू. $\pi$ मधील अस्वच्छ आयगेन
चराच्या निगमनांची संख्या~$n$ असो. $n$ वर
विगमन करून~$\pi$ चे अपेक्षित नियमित
!!{proof}~$\pi'$ मध्ये रूपांतर करता येते हे
दाखवू. $n=0$ असेल तर~$\pi$ मधील आयगेन चराची
सर्व निगमने स्वच्छ आहेत; म्हणून~$\pi$ नियमित
असते आणि सिद्ध करण्यास काही उरत नाही.
अन्यथा, सर्वांत वरचे अस्वच्छ आयगेन-चर निगमन,
उदा.~\RightR{\lforall}, निवडा. त्यात संपणारी
~$\pi$ ची उप-!!{proof}~$\pi_1$ पुढील रूपाची आहे:
\[
    \AxiomC{}
    \RightLabel{$\pi_2$}
    \Deduce$\Gamma \fCenter \Delta, !A(c)$
    \RightLabel{\RightR{\lforall}}
    \UnaryInf$\Gamma \fCenter \Delta, \lforall[x][!A(x)]$
    \DisplayProof
    \]
$c$ हा आयगेन चर असल्यामुळे तो~$\Gamma$,
$\Delta$, किंवा $\lforall[x][!A(x)]$ मध्ये
येत नाही. $\pi_2$ नियमित आहे, कारण तिच्यातील
कोणतेही आयगेन-चर निगमन अस्वच्छ नाही.
~$\pi$ मध्ये न येणारा !!a{constant}~$a$ घ्या.
\olref{lem:subst-regular} नुसार
$\Subst{\pi_2}{a}{c}$ ही $\Gamma \fCenter \Delta,
!A(a)$ ची नियमित सिद्धता आहे; तिच्यावर
~$\RightR{\lforall}$ लावता येतो. $\pi$ मधील
~$\pi_1$ च्या जागी पुढील सिद्धता~$\pi_1'$ ठेवू:
\[
    \AxiomC{}
    \RightLabel{$\Subst{\pi_2}{a}{c}$}
    \Deduce$\Gamma \fCenter \Delta, !A(a)$
    \RightLabel{\RightR{\lforall}}
    \UnaryInf$\Gamma \fCenter \Delta, \lforall[x][!A(x)]$
    \DisplayProof
    \]
त्यामुळे एक अस्वच्छ आयगेन-चर निगमन स्वच्छ बनते;
फरक फक्त आयगेन चर~$c$ च्या जागी~$a$ आल्याचा
आहे. परिणामी सिद्धतेतील अस्वच्छ आयगेन-चर
निगमनांची संख्या~$n-1$ पेक्षा जास्त राहत नाही,
म्हणून विगमनात्मक गृहीतकातून निष्कर्ष मिळतो.
\end{proof}

नुकताच सिद्ध केलेला निष्कर्ष आपण स्पष्टपणे
संदर्भ म्हणून वापरणार नाही; पण पुढील सर्व
!!{proof}s नियमित आहेत असे मानू. त्या नियमित
नसल्या तर आयगेन चर बदलून त्यांना नेहमी नियमित
बनवता येते.

\end{document}
